\documentclass[10pt,twocolumn]{article}
\usepackage[margin=0.72in]{geometry}
\usepackage{booktabs}
\usepackage{microtype}
\usepackage{url}
\usepackage{xcolor}
\usepackage{graphicx}
\usepackage{amsmath}

\title{Shadow Trainer: Auditable, Resource-Bounded Training Across GPU Memory, Host Memory, Disk, and Remote Storage}
\author{Anonymous submission}
\date{}

\begin{document}
\maketitle

\begin{abstract}
Training pipelines commonly assume that a dataset is locally available and
that failures can be handled independently from memory management. These
assumptions break on edge workstations, laboratory computers, and development
laptops where GPU memory, host memory, disk capacity, and network data locality
are all hard constraints. We present Shadow Trainer, a local runtime that
admits a job only when its declared resource contract is safe, stages complete
cases through an integrity-checked bounded cache, creates durable checkpoints
that include optimizer and random-number-generator state, and emits a
machine-verifiable execution record. A conservative policy keeps overlapped
prefetch experimental until a paired equivalence gate passes. We report an
initial implementation and physical validation on a laptop NVIDIA RTX 3050 Ti
with 4 GiB of VRAM. A portable 3D CNN fixture completed 12 steps while respecting
a 4 KiB cache, and a bounded NIfTI workload completed two real 3D image-label
cases using $64^3$ patches. ResNet-18 and ResNet-50 topology adapters also
completed short 2.5D proxy calibrations through the same public runtime. Peak reserved GPU memory was 22 and 64 MiB for the 3D runs, and 258 and 488 MiB for the ResNet calibrations. These measurements establish execution
feasibility and artifact completeness, not throughput superiority or clinical
quality. We also preserve a deterministic-CUDA failure and an incomplete
prefetch experiment as negative evidence. The resulting system and evaluation
protocol separate product readiness from optimization claims and provide a
foundation for broader evaluation across models and storage conditions.
\end{abstract}

\section{Introduction}

Commodity accelerators make local machine learning attractive for privacy,
cost, iteration speed, and access. Their nominal compute capability, however,
does not make a training job operationally safe. A workstation can have enough
free GPU memory for a patch-based model yet insufficient disk space for a
dataset replica, or sufficient disk capacity but insufficient host memory and
swap headroom during decoding. Remote datasets add partial downloads,
credentials, transient connectivity, and cache eviction to the failure model.
Restarting after a failure can silently repeat samples or change the random
sequence unless model state, optimizer state, data position, and random state
are treated as one recovery unit.

Existing systems address important slices of this problem. ZeRO-Offload moves
model state and computation between GPU and CPU to enable much larger language
models~\cite{zerooffload}. CheckFreq provides fine-grained DNN checkpointing
and a resumable data iterator~\cite{checkfreq}. MONAI offers domain-specific
medical-imaging pipelines and caching abstractions~\cite{monai}. Recent work on
edge retraining streams optimizer state from storage to reduce unified-memory
pressure~\cite{steer}. Shadow Trainer explores a different systems boundary:
the admission, data-locality, recovery, and audit contract for a complete local
job whose constraints span GPU memory, RAM, swap, disk, and a remote source.

This paper makes four preliminary contributions:

\begin{enumerate}
  \item A versioned job contract and deterministic admission policy that reject
  unsafe jobs before data transfer or GPU allocation.
  \item An atomic, hash-validated case stager with a byte-accurate LRU cache and
  explicit synchronous, prefetch, and conservative automatic modes.
  \item A durable case-window recovery boundary that captures model, optimizer,
  RNG, and dataset position together.
  \item An evidence compiler that validates run bundles, derives metrics from
  event logs, hashes their provenance, and excludes diagnostic or invalid runs
  from numeric figures.
\end{enumerate}

Our initial experiments answer a deliberately narrow question: can the runtime
enforce its basic contracts on physical 4-GiB hardware? They do not yet
establish speedup, scale independence, full-state equivalence under prefetch,
or clinical value.

\section{Problem and threat model}

We consider a single Linux or WSL host with one NVIDIA GPU. A job consumes a
manifest of cases from either a local directory or an \texttt{rclone}-addressed
source. The complete dataset may be larger than the allowed cache and available
local disk. Training is organized into case windows, but the workload adapter
controls model-specific tensor construction.

The runtime must tolerate interruption at a committed window boundary,
transient source failure, truncated transfer, hash mismatch, and resource
shortage detected before execution. It must not promote a partial file into the
cache, exceed the declared cache or combined-artifact budget, traverse outside
the cache root, expose source credentials in reports, or claim a step that was
not durably checkpointed.

We do not defend against a malicious kernel, compromised Python process,
incorrect manifest supplied by a trusted operator, or adversarial training
code inside a workload adapter. Remote availability is not guaranteed. The
current implementation also does not resize tensors or models automatically;
it admits or rejects a declared workload and orchestrates its data and state.

\section{Design}

\subsection{Versioned resource contract}

A \texttt{shadowtrainer.job/v1} document identifies the job and seed, output
directory, data manifest and source, cache budget, disk floor, minimum free RAM
and GPU memory, maximum swap use, artifact ceiling, execution strategy,
training windows, and workload adapter. Byte counts are integers to avoid
ambiguous unit parsing. Before creating a run directory, the planner captures
GPU, driver, CUDA, cuDNN, RAM, swap, disk, Python, and dependency state. A
rejected plan returns structured reasons without staging a case or constructing
the model.

The \texttt{auto} strategy is intentionally asymmetric. It selects synchronous
staging while the equivalence gate for prefetch remains open. Users can request
prefetch for controlled experiments, but the plan marks it experimental. This
turns scientific uncertainty into executable product policy rather than a
documentation caveat.

\subsection{Atomic case staging}

The manifest maps stable case identifiers to named files, sizes, and SHA-256
digests. A local or remote source copies each file to a temporary path. Only a
file with the expected size and digest is atomically renamed into the cache.
The stager validates that resolved paths remain below the configured root and
serializes cache mutation. LRU eviction considers complete cases rather than
arbitrary file fragments. The current and next windows may be protected from
eviction only if their combined declared size fits the cache.

Synchronous mode stages the current window on demand. Prefetch mode uses one
background worker for the next window and records launch, consumption, or skip
events. It skips overlap when protecting both windows would exceed the byte
budget. This policy prioritizes bounded behavior over nominal overlap.

\subsection{Durable execution and recovery}

After each case window, the runtime atomically writes a checkpoint containing a
digest of the public job configuration, the next epoch and window, the global
step, and workload-owned state. The provided PyTorch adapters include model,
optimizer, Python RNG, CPU Torch RNG, and all CUDA RNG states. A resume operation
requires a matching configuration digest. The event stream appends a checkpoint
record only after the atomic save completes. The runtime then checks combined
cache and artifact bytes against the declared ceiling.

This boundary is coarser than CheckFreq's adaptive iteration checkpointing but
also includes case-locality state. The design can later incorporate an adaptive
frequency policy without changing the public recovery contract.

\subsection{Workload adapters}

The interface consists of \texttt{train\_case}, \texttt{checkpoint\_state}, and
\texttt{restore}. Tiny3D is a deterministic portable CNN used for offline smoke
tests. The NIfTI adapter memory-maps image and label volumes, selects a
deterministic bounded patch, and materializes only that patch on the accelerator.
Its compact four-class CNN is a systems probe rather than a competitive
segmentation architecture.

An initial CUDA path using a standard cross-entropy operation failed because
the installed deterministic backward implementation was unavailable. We
preserved that attempt and used the equivalent negative log-probability
formulation based on \texttt{log\_softmax} and \texttt{gather} in a new run.
This illustrates why failure evidence and environment capture are part of the
runtime rather than post-hoc experiment notes.

\subsection{Evidence compiler}

Every run emits the public job, plan, environment, JSONL events, summary,
checkpoint, and static HTML report. The evidence compiler requires the first
five files, parses every event, derives metrics, and hashes names and contents
into one SHA-256 provenance digest. A separate specification classifies each
run as valid, engineering-only, diagnostic, invalid, or pending and states its
permitted use. Numeric SVG figures include only valid records. Paths to source
data and caches are not written to generated tables.

\section{Correctness invariants}

Shadow Trainer is intended to make operational guarantees testable. The
implementation is therefore organized around invariants rather than a target
throughput number.

\paragraph{Admission non-interference.}
A rejected job must not materialize source objects, instantiate its workload,
or allocate accelerator state. The planner operates only on parsed configuration
and a resource snapshot. This ordering matters when admission protects a nearly
full disk: probing feasibility by first creating a cache would itself consume
the resource being protected. The physical low-disk observation in this study
returned \texttt{disk\_headroom} before staging. Unit tests separately force
GPU and disk rejection paths.

\paragraph{Atomic promotion.}
A source object becomes cache-visible only after transfer to a temporary name,
size validation, optional format validation, digest verification, and atomic
rename. A timeout, process exit, truncated object, or mismatched digest can
leave reconstructible temporary data but cannot produce a valid cache entry.
The cache state records complete cases, and a hit is accepted only for files
that still meet the manifest contract. This is stronger than assuming that a
successful copy command implies a usable training sample.

\paragraph{Bounded residency.}
Before admitting a new case, the stager computes the required bytes and evicts
least-recently-used unprotected cases until the postcondition fits the declared
budget. A case larger than the entire cache is rejected rather than partially
resident. During prefetch the current and future windows are both protected;
overlap is skipped when their declared sum exceeds the cache. At a durable
boundary the runtime also measures the checkpoint and all run artifacts together
with cache occupancy and compares the result with the job-level artifact ceiling.

The current implementation checks this combined ceiling at window boundaries,
not continuously while a large checkpoint temporary file is being written.
Atomic save protects checkpoint validity, but a future version should reserve
the predicted temporary-file size before writing so that instantaneous disk
occupancy is covered by the same hard invariant.

\paragraph{Monotone committed progress.}
The checkpoint stores the next epoch and window, not merely the most recent
global step. On resume, a configuration digest prevents accidental continuation
under a different manifest, resource contract, workload, or output identity.
The event log can contain multiple process-start records, but confirmed training
steps must remain monotone and unique. The paired Tiny3D test now compares final
model tensors, optimizer state, Python RNG state, and Torch RNG state between an
uninterrupted execution and a one-step interruption followed by resume. Equality
is established for this adapter only; the same proof obligation remains for
NIfTI and ResNet.

\paragraph{Evidence non-escalation.}
Completing a process is not sufficient to authorize every statement about it.
Each evidence record has a classification and permitted use independent of the
run's status. A successful two-case medical-image run can support systems
feasibility but not clinical quality. A failed run may explain a design change
but cannot enter a numeric figure. A complete sync arm without its paired
prefetch arm cannot support a speedup. The compiler enforces the simplest of
these rules mechanically by selecting only \texttt{valid} records for SVG
figures; the claims ledger constrains interpretation further.

These invariants expose remaining proof gaps. The file and cache layer has
failure-injection tests, whereas remote recovery across an actual connectivity
loss is pending. The runtime restores workload state exactly for Tiny3D, whereas
data transforms and asynchronous workers may add nondeterminism in larger
adapters. Finally, the resource snapshot is a point observation. An external
process can consume memory or disk after admission, so long jobs also require
runtime guards and a clearly reported guard stop rather than relying exclusively
on the initial plan.


\section{Evaluation methodology}

We ask whether admission is safe (RQ1), cache occupancy is bounded (RQ2),
recovery is idempotent (RQ3), prefetch preserves complete observable behavior
(RQ4), and the contract generalizes across workloads (RQ5). The current
evaluation addresses the first three partially and provides initial multi-workload
support for the fifth.

The physical host runs WSL2 with an NVIDIA RTX 3050 Ti Laptop GPU (4 GiB),
driver 610.62, PyTorch 2.12.0 built for CUDA 13.0, and cuDNN 92000. P1 is a
portable six-case fixture trained for two epochs. Its four-case cache forces
eviction. P2 contains two KiTS23 NIfTI image-label pairs and uses one $64^3$
patch per case, mixed precision, and a 128-MiB cache. Dataset content is not
reported. D1 is the preserved failed form of P2 and is diagnostic only.

Table~\ref{tab:initial} reports values derived from immutable JSON/JSONL
artifacts. The bundle digests are generated alongside CSV, Markdown, and SVG
outputs. Since P1 and P2 use different data and workloads, their times and
memory values must not be interpreted as a comparison.

\begin{table*}[t]
\centering
\caption{Initial physical evidence. All rows are independent feasibility runs, not comparative arms.}
\label{tab:initial}
\begin{tabular}{llrrrrrr}
\toprule
ID & Workload & Steps & Time (s) & Peak alloc. (MiB) & Peak reserved (MiB) & Cache used/budget & Transfer (MiB) \\
\midrule
P1 & Tiny3D & 12/12 & 2.316 & 16.48 & 22.00 & 4096/4096 B & 0.012 \\
P2 & NIfTI patch 3D & 2/2 & 4.816 & 47.60 & 64.00 & 79.68/128 MiB & 36.68 \\
P3 & ResNet-18 2.5D proxy & 4-step limit & 6.813 & 235.15 & 258.00 & 8/8 KiB & 0.008 \\
P4 & ResNet-50 2.5D proxy & 2-step limit & 5.278 & 471.00 & 488.00 & 4/8 KiB & 0.004 \\
\bottomrule
\end{tabular}
\end{table*}

All four valid runs completed their configured full or step-limited executions with finite losses and generated
the required report and checkpoint artifacts. P1 reached exactly its cache
budget and performed repeated whole-case eviction without exceeding it. P2 had
one warm case and one cache miss; its measured staging time was 1.076 seconds.
The evidence is sufficient for physical feasibility and bounded-cache claims.
It is insufficient for throughput, capacity beyond local storage, or clinical
segmentation claims.

Unit and integration tests cover strict configuration parsing, manifest and
path validation, atomic partial-transfer behavior, LRU enforcement, admission
rejection, rclone success and timeout behavior, checkpoint resume without
repeating confirmed step identifiers, report redaction, the NIfTI adapter, and
evidence bundle integrity. The current suite contains 25 tests. A paired Tiny3D test establishes exact model, optimizer, and RNG equality after interruption and resume. Small CPU fixtures also pass exact sync/prefetch audits for all four workload configurations.

\subsection{Open prefetch gate}

A historical paired protocol prepared independent caches and launched a
synchronous arm intended to precede prefetch. The sync process completed 108 of
144 steps and resumed at step 72, then a supervisor terminated it when physical
disk free space fell below the 20-GiB floor. No OOM, non-finite value, or remote
error was observed. The prefetch arm was not launched. Consequently there is
no valid pair and no speedup or equivalence result. Production \texttt{auto}
therefore remains synchronous.

\section{Protocol for comparative evaluation}

A performance evaluation must separate three effects that are easily conflated:
capacity, runtime overhead, and latency hiding. Shadow Trainer may enable a job
under a disk ceiling even when it is slower than framework-local loading. A
prefetch arm may reduce exposed I/O time while leaving total runtime unchanged
because training dominates. Conversely, a warm source or filesystem cache may
appear to accelerate prefetch even when overlap provides no benefit. We define
the following protocol to prevent these cases from becoming one aggregate
``speedup'' number.

\subsection{Arms and controlled variables}

The framework-local arm reads a fully local dataset without Shadow Trainer and
establishes orchestration overhead. The synchronous arm uses atomic staging and
the bounded cache. The prefetch arm changes only whether the next case window is
materialized concurrently. Paired sync and prefetch processes use the same
manifest bytes, case order, seed, initial checkpoint, model options, training
steps, cache capacity, disk floor, driver, CUDA, cuDNN, and PyTorch version.
Each arm starts in a new process to avoid allocator and library state leaking
between measurements.

Cold and warm conditions are separate populations. A cold run starts from an
independently prepared cache state and records every source byte. A warm run is
allowed to reuse the declared case cache but must identify its initial contents.
Operating-system page-cache state is difficult to control safely on a shared
workstation; it is therefore reported as a limitation rather than manipulated
with privileged cache-dropping commands.

\subsection{Measurements}

Primary capacity outcomes are admission, completion, peak allocated and reserved
VRAM, process RSS, swap use, maximum cache occupancy, combined artifact bytes,
and minimum observed disk headroom. Primary recovery outcomes are lost or
repeated step identifiers and equality of final model, optimizer, RNG, and data
position. Performance outcomes are end-to-end time, staging time, bytes moved,
cache hits and misses, exposed I/O stall, steps per second, and GPU utilization.

Every reported performance cell requires at least three independent processes;
five are preferred for short workloads. We report individual observations,
median, interquartile range, and a bootstrap confidence interval for the paired
difference. A mean alone is not used for small, skewed I/O samples. Hardware
telemetry sampling must be identical across arms and sufficiently light that its
overhead is either measured or held constant.

\subsection{Exactness gate}

Latency hiding is eligible for a performance comparison only after the complete
observable result passes an exactness audit. The audit checks ordered case and
step identifiers, finite losses, final model and optimizer tensors, RNG states,
checkpoint position, manifest digest, and boundary hashes. If deterministic
algorithms cannot make a particular kernel bitwise stable, the protocol must
predeclare numerical tolerances and justify them; a tolerance cannot be selected
after observing the difference. The current product takes the stricter position:
prefetch remains experimental because the full pair has not completed.

\subsection{Capacity experiment}

The principal paper result should use a remote dataset whose declared size
exceeds the cache and the disk allocation available to the job. Capacity is
demonstrated only if all requested cases are processed, maximum cache and
artifact occupancy remain within contract, every promoted file passes integrity
validation, and the final disk floor is respected. Comparing total dataset size
with physical disk size is unnecessary and potentially misleading; the relevant
quantity is the storage allocation made available under the experimental
contract.

The same experiment should be repeated for NIfTI 3D, ResNet-18, ResNet-50, and
a real adjacent-slice 2.5D pipeline. The current proxy ResNet runs validate the
adapter boundary and topology memory footprint but do not replace real decoding,
augmentation, or data-loader behavior. Together, the capacity and overhead arms
would determine whether the unified contract provides a useful capability and
what performance price it imposes.


\section{Related work}

Memory-efficient model training often focuses on model state. ZeRO-Offload
partitions and offloads optimizer computation and state to CPU memory while
minimizing GPU--CPU movement~\cite{zerooffload}. Out-of-core scheduling also
moves model variables between GPU and CPU when they are not needed
immediately~\cite{outofcore}. STEER targets edge retraining and streams
optimizer state from fast storage~\cite{steer}. Shadow Trainer does not yet
offer comparable model-state capacity. Its focus is a job-level contract across
resources and dataset locality, so these techniques are complementary candidates
for future workload adapters.

Checkpointing research addresses recovery cost and data semantics. CheckFreq
profiles checkpoint overhead, adapts frequency, pipelines persistence, and
restores data-loader position~\cite{checkfreq}. Shadow Trainer currently uses a
fixed case-window boundary, adds source/cache integrity and admission evidence,
and treats the resulting bundle as an auditable product artifact.

Medical-imaging frameworks such as MONAI provide transformations, architectures,
and cache-oriented datasets on PyTorch~\cite{monai}. Shadow Trainer uses NIfTI
as one physical validation workload but is domain-neutral. A fair evaluation
must compare framework-local MONAI loading with Shadow Trainer's bounded remote
staging rather than presenting the frameworks as substitutes.

\section{Limitations and next experiments}

The study has one hardware host, four successful short runs, and no statistically
repeated performance experiment. The compact NIfTI model does not establish
medical utility. GPU memory values are workload observations, not maximum
supported capacity. The current cache trace does not yet demonstrate a remote
dataset larger than all available local storage. Recovery tests prove step-id idempotence and bitwise final model, optimizer, and
RNG equality for Tiny3D, but a general claim requires equivalent paired tests
for the NIfTI and ResNet adapters. Prefetch has not passed its gate.

The next evaluation uses separate processes and independently prepared caches
for sync and prefetch, with identical manifest, order, seed, checkpoint, CUDA,
cuDNN, and driver. Cold and warm runs remain separate. The runtime will stop
before preparation unless physical disk has at least 30 GiB free, retain a
20-GiB floor, and cap combined cache and artifacts at 5 GiB. ResNet18, ResNet50,
and a 2.5D adapter will test generality. A framework-local PyTorch/MONAI arm
will quantify overhead. Only a pair with identical observable state can support
an overlap-performance claim.

\section{Conclusion}

Shadow Trainer makes resource limits, dataset locality, recovery, and evidence
validity explicit parts of a local training job. Initial 4-GiB GPU runs show
that the implementation can execute a synthetic 3D CNN, bounded real NIfTI
patch training, and ResNet-18/50 2.5D proxy calibrations while producing complete, traceable artifacts. More importantly,
the same policy prevents an unfinished prefetch experiment from becoming a
product speedup claim. The remaining work is empirical: establish paired
equivalence, evaluate larger-than-cache remote data, add workload families, and
measure overhead across repeated hardware runs.

\bibliographystyle{plain}
\bibliography{references}
\end{document}
